\documentclass[11pt,a4paper]{article}

% --- Packages ---
\usepackage[utf8]{inputenc}
\usepackage[margin=2.2cm]{geometry}
\usepackage{booktabs}
\usepackage{graphicx}
\usepackage{amsmath}
\usepackage{amssymb}
\usepackage{hyperref}
\usepackage{xcolor}
\usepackage{microtype}
\usepackage{caption}
\usepackage{fancyhdr}

% --- Hyperlink Setup ---
\hypersetup{
    colorlinks=true,
    linkcolor=blue!70!black,
    citecolor=blue!70!black,
    urlcolor=blue!70!black
}

% --- Page Setup ---
\pagestyle{fancy}
\fancyhf{}
\rhead{\small KIUA2003 — Compulsory Assignment 1}
\lhead{\small Multi-Agent Interrogation System}
\cfoot{\thepage}

\title{\textbf{Compulsory Assignment 1: Multi-Agent Dialogue System}\\[0.4ex]
\large Building, Evaluating, and Breaking a Two-Agent Interrogation Loop}
\author{\textbf{Course:} KIUA2003 — Applikasjon av kunstig intelligens / maskinl\ae{}ring II}
\date{September 2026}

\begin{document}

\maketitle

% ==============================================================================
% SECTION 1: INTRODUCTION
% ==============================================================================
\section{Introduction}

In this assignment, we built and tested a two-agent dialogue system from scratch in Python without relying on pre-built agent frameworks like AutoGen or LangGraph. The main goal was to understand how multi-agent interaction works under the hood: how agents alternate turns, how conversation history is formatted for each role, how budget limits prevent runaway loops, and how agents behave when pushed into common failure modes.

We chose a criminal interrogation scenario between a detective and a suspect. This creates natural conflict and clear goals: the detective tries to find holes in an alibi, while the suspect tries to defend their story without confessing. In this report, we describe our system design, present experimental results across different temperature settings, analyze three deliberate failure modes, examine token and time costs, and reflect on what existing frameworks automate for developers.

% ==============================================================================
% SECTION 2: SCENARIO AND SYSTEM DESIGN
% ==============================================================================
\section{Scenario and System Design}

\subsection{The Interrogation Scenario}
Our scenario centers on the theft of a priceless diamond, the \textit{``Star of Midnight''}, from a museum vault at 21:30. The two agents are:
\begin{itemize}
    \item \textbf{Detective Cross (Interrogator):} A skeptical investigator asking focused questions to test timeline inconsistencies.
    \item \textbf{Julian Vance (Prime Suspect):} The chief museum curator who claims he was dining alone across town at the Grand Bistro between 21:00 and 22:30.
\end{itemize}

\begin{figure}[htbp]
    \centering
    \IfFileExists{diagram.png}{%
        \includegraphics[width=0.85\linewidth]{diagram.png}%
    }{%
        \fbox{
            \begin{minipage}{0.88\linewidth}
                \centering
                \vspace{0.8cm}
                \textbf{\Large [ Scenario Diagram: \texttt{diagram.png} ]} \\[0.3cm]
                \small \textit{Upload \texttt{diagram.png} to Overleaf to render this figure automatically.}
                \vspace{0.8cm}
            \end{minipage}
        }%
    }
    \caption{Overview of the Diamond Theft Interrogation Scenario and Agent Roles.}
    \label{fig:scenario_diagram}
\end{figure}

\subsection{Why We Chose This Scenario}
We selected an interrogation rather than a symmetric debate for three reasons:
\begin{enumerate}
    \item \textbf{Natural Conflict and Information Asymmetry:} In symmetric debates, agents often end up agreeing too quickly. In an interrogation, the suspect has private knowledge about their whereabouts, while the detective has institutional authority and forensic clues.
    \item \textbf{Clear Conversational States:} Progress is easy to observe and measure: the suspect either maintains their alibi, contradicts themselves, changes their story, or gets backed into a corner.
    \item \textbf{Resistance to Quick Agreement:} Because the agents have directly opposing incentives, the dialogue provides a good test of whether models stay in character instead of falling into polite agreement.
\end{enumerate}

\subsection{Agent Prompts and Settings}
Each agent has a system prompt, a model identifier (\texttt{llama3.2:3b}), and a temperature setting.

\subsubsection*{Detective Cross (Interrogator)}
\begin{itemize}
    \item \textbf{System Prompt:}
    \begin{quote}
    \small \textit{``You are Detective Cross interrogating Julian Vance about the theft of the 'Star of Midnight' diamond from the Grand Gallery vault at 21:30. Ask direct, pointed questions probing his whereabouts and uncover any inconsistencies in his alibi. Be concise: 1 to 2 sentences per response.''}
    \end{quote}
\end{itemize}

\subsubsection*{Julian Vance (Prime Suspect)}
\begin{itemize}
    \item \textbf{System Prompt:}
    \begin{quote}
    \small \textit{``You are Julian Vance, chief curator and prime suspect in the theft of the 'Star of Midnight' diamond. Your alibi is that you were having dinner alone across town at the Grand Bistro between 21:00 and 22:30. Defend your alibi, respond cautiously, and do not admit guilt unless backed into a corner. Be concise: 1 to 2 sentences per response. At the very end of every reply, append a JSON block: \{"confessed": false\} (or true if you admit guilt)''}
    \end{quote}
\end{itemize}

\subsection{How the Dialogue Engine and Budgets Work}
The core loop in \texttt{engine.py} combines four main components:
\begin{enumerate}
    \item \textbf{The Mirror Property (\texttt{view\_for}):} Chat models require messages labeled with standard roles (\texttt{system}, \texttt{user}, \texttt{assistant}). The \texttt{view\_for} function maps the shared conversation so that an agent's own past replies become \texttt{assistant}, and the opposing agent's replies become \texttt{user}.
    \item \textbf{Context Management (\texttt{manage\_context}):} If a dialogue continues for many turns, the prompt eventually exceeds the context window. We implemented \texttt{summarise\_context}, which compresses older messages into a short summary at temperature 0 while keeping recent turns intact.
    \item \textbf{Budget Guardrails (\texttt{Budget}):} To prevent infinite loops or excessive compute use, the \texttt{Budget} class tracks three hard limits: turns ($\le 8$), tokens ($\le 4000$), and elapsed wall-clock seconds ($\le 120$). It updates incrementally after each turn.
    \item \textbf{Goal Check and Parse Guard:} We check for a confession using regular expressions and a \texttt{json.loads()} try-catch block on the suspect's structured JSON field. This prevents premature stops that would happen with simple keyword matching (such as stopping if the suspect says \textit{``I will never confess''}).
\end{enumerate}

% ==============================================================================
% SECTION 3: RESULTS AND DISCUSSION
% ==============================================================================
\section{Results and Observations}

\subsection{Overview of Test Runs}
All results in Table~\ref{tab:results} were extracted directly from our saved JSON transcript files in \texttt{transcripts/} using the local \texttt{llama3.2:3b} model. We tested our baseline configuration, a temperature sweep from 0.0 to 1.0, and three deliberate failure modes.

\begin{table*}[htbp]
\centering
\small
\caption{Results Across Experimental and Failure Configurations.}
\label{tab:results}
\resizebox{\linewidth}{!}{%
\begin{tabular}{llccccccccc}
\toprule
\textbf{Configuration} & \textbf{Description} & \textbf{Temp ($T$)} & \textbf{Turns} & \textbf{Prompt} & \textbf{Comp.} & \textbf{Total} & \textbf{Time (s)} & \textbf{Judge} & \textbf{Success} & \textbf{Stop Reason} \\
\midrule
\multicolumn{11}{l}{\textit{\textbf{Baseline and Temperature Tests}}} \\
\textbf{\texttt{final.yaml}} & \textbf{Agreed Baseline} & \textbf{0.3 / 0.7} & \textbf{8} & \textbf{3,087} & \textbf{591} & \textbf{3,678} & \textbf{61.14} & \textbf{2 / 5} & \textbf{False} & \textbf{\texttt{max\_turns}} \\
\texttt{exp-temp00.yaml} & Low Temp (Greedy) & 0.0 & 8 & 3,222 & 625 & 3,847 & 24.52 & 2 / 5 & False & \texttt{max\_turns} \\
\texttt{exp-temp07.yaml} & Medium Temp & 0.7 & 8 & 2,968 & 602 & 3,570 & 24.14 & 2 / 5 & False & \texttt{max\_turns} \\
\texttt{exp-temp10.yaml} & High Temp & 1.0 & 8 & 2,940 & 607 & 3,547 & 24.27 & \textbf{4 / 5} & \textbf{True} & \texttt{max\_turns} \\
\midrule
\multicolumn{11}{l}{\textit{\textbf{Failure Mode Benchmarks}}} \\
\texttt{fail-drift.yaml} & Extreme Temp (Drift) & 1.5 & 8 & 2,734 & 574 & 3,308 & 23.83 & 2 / 5 & False & \texttt{max\_turns} \\
\texttt{fail-forgetting.yaml} & Truncation (30T limit) & 0.3/0.7 & 12 & 3,352 & 784 & 4,136 & 36.63 & 2 / 5 & False & \texttt{max\_tokens} \\
\texttt{fail-sycophancy.yaml} & Overly Agreeable & 0.3/0.7 & 8 & 2,447 & 465 & 2,912 & 24.52 & 4 / 5$^*$ & True$^*$ & \texttt{max\_turns} \\
\bottomrule
\end{tabular}%
}
\begin{flushleft}
\footnotesize $^*$Note: The 4/5 score on \texttt{fail-sycophancy} illustrates an evaluator blindspot, explained in Section~\ref{sec:failures}.
\end{flushleft}
\end{table*}

\subsection{What the Results Show}
\begin{enumerate}
    \item \textbf{Baseline Behavior (\texttt{final.yaml}):} Setting the detective at $T=0.3$ (calm, methodical) and the suspect at $T=0.7$ (creative defense) produced a realistic interrogation. Detective Cross systematically pressed Vance on phone records, car parking footage, a fake repair shop alibi, and a sudden GPS slowdown at 21:30. Vance was forced through multiple alibi changes and was hesitating noticeably by Turn 8. Because he did not explicitly confess, the judge gave a 2/5 score, but noted that the suspect's alibi was completely broken.
    \item \textbf{Temperature and Repetition Loops:} At $T=0.0$, the small 3B model got stuck repeating the exact same phrase (\textit{``I must have misremembered...''}) across four consecutive turns. Increasing the temperature to 0.7 and 1.0 gave the model enough variety to avoid loops and introduce plausible details (such as names of museum colleagues), scoring higher with the judge.
    \item \textbf{Prompt Tokens Dominate Cost:} In all 8-turn dialogues, prompt tokens made up roughly $84\%$ of total token usage. This happens because the whole conversation history is re-sent on every new turn, making total prompt tokens grow rapidly over time while generated replies stay around 75 tokens per turn.
    \item \textbf{Budget Enforcement Under Stress:} In \texttt{fail-forgetting.yaml}, the run was configured for up to 30 turns. However, the dialogue stopped automatically at Turn 12 when it hit 4,136 tokens, showing that our token limit guardrail works as intended.
\end{enumerate}

% ==============================================================================
% SECTION 4: FAILURE-MODE ANALYSIS
% ==============================================================================
\section{Failure Mode Analysis}
\label{sec:failures}

As highlighted in the Week 4 brief, identifying and analyzing failure modes is a key part of understanding multi-agent systems. We deliberately drove our system into three common failures:

\subsection{1. Topic and Persona Drift (\texttt{fail-drift.json})}
\begin{itemize}
    \item \textbf{What Happened:} At an extreme temperature of $T=1.5$, Julian Vance completely abandoned his Grand Bistro dinner alibi in Turn 1:
    \begin{quote}
    \small \textit{``I was attending a conference at the Grand Gallery, actually presenting a paper on gemology, and I left the venue around 21:15.''} \hfill --- [Turn 1]
    \end{quote}
    By Turn 7, he claimed he was rushing to catch an early flight.
    \item \textbf{Cause:} Extremely high temperature makes next-token selection nearly random. The model stops paying attention to the system prompt instructions and drifts away from the topic.
    \item \textbf{Mitigation:} Keep temperature at or below 0.7, and restate the key goal or alibi constraint in the user prompt each turn.
\end{itemize}

\subsection{2. Forgetting After Context Compression (\texttt{fail-forgetting.json})}
\begin{itemize}
    \item \textbf{What Happened:} In an extended 12-turn run with context compression, Vance began contradicting his earlier alibi (shifting arrival times and admitting to wearing a gold watch that belonged to a bystander).
    \item \textbf{Cause:} Summarization is lossy. When older messages are compressed, specific details like exact minutes and alibi claims are dropped, leading to inconsistencies.
    \item \textbf{Mitigation:} Instead of relying solely on text summaries, maintain a structured list of established facts (such as a simple JSON ledger of times and locations).
\end{itemize}

\subsection{3. Sycophancy and Judge Blindspot (\texttt{fail-sycophancy.json})}
\begin{itemize}
    \item \textbf{What Happened:} When we gave both agents agreeable system prompts, Detective Cross opened politely: \textit{``Good morning Mr.~Vance... I appreciate your cooperation.''} Vance immediately replied: \textit{``It's an absolute pleasure to help...''}. The interrogation collapsed into mutual compliments and paperwork, with zero investigation taking place.
    \item \textbf{The Evaluator Mistake:} Surprisingly, our automated LLM judge awarded this run a \textbf{4/5 (Success)}, stating that the detective conducted a \textit{``coherent and probing inquiry''}.
    \item \textbf{Cause:} Chat models are trained via RLHF to be polite and helpful. The judge model was misled by the polite tone and failed to notice that no actual interrogation happened.
    \item \textbf{Mitigation:} Add explicit negative scoring criteria to judge prompts, such as requiring evidence of challenged statements before awarding a passing score.
\end{itemize}

% ==============================================================================
% SECTION 5: COST AND TIME ACCOUNTING
% ==============================================================================
\section{Cost and Time Accounting}

Running local models is free in monetary terms, but it consumes time and compute resources. On our local setup running \texttt{llama3.2:3b}, each turn took roughly 3.0 seconds.

\subsection{Estimating Costs on Paid APIs}
To understand what this workload would cost on commercial APIs, we use the standard token pricing formula:
\begin{equation}
\text{Cost} = \left(\frac{\text{Prompt Tokens}}{1000} \times \text{Price}_{\text{prompt}}\right) + \left(\frac{\text{Completion Tokens}}{1000} \times \text{Price}_{\text{completion}}\right)
\end{equation}
Using typical commercial rates (GPT-4o at $\$0.005$ per 1K prompt and $\$0.015$ per 1K completion; GPT-3.5-Turbo at $\$0.0005$ prompt and $\$0.0015$ completion), our average 8-turn run costs:
\begin{itemize}
    \item \textbf{GPT-4o Estimate:} $\approx \$0.015 + \$0.009 = \mathbf{\$0.024 \text{ per run}}$.
    \item \textbf{GPT-3.5-Turbo Estimate:} $\approx \$0.0015 + \$0.0009 = \mathbf{\$0.0024 \text{ per run}}$.
\end{itemize}
\textbf{Key Takeaway:} Because the conversation history grows with every turn, prompt tokens account for the majority of the cost ($>60\%$). In multi-agent applications, managing the context window is essential to avoid ballooning API bills.

% ==============================================================================
% SECTION 6: REFLECTION AND FRAMEWORK COMPARISON
% ==============================================================================
\section{Reflection and Framework Comparison}

Building our own orchestration loop made it clear what modern multi-agent frameworks handle behind the scenes. Table~\ref{tab:framework_mapping} maps our handcrafted Python components to their equivalents in AutoGen, LangGraph, and CrewAI.

\begin{table}[htbp]
\centering
\small
\caption{Mapping Our Python Code to Multi-Agent Frameworks.}
\label{tab:framework_mapping}
\resizebox{\linewidth}{!}{%
\begin{tabular}{llll}
\toprule
\textbf{Our Code Component} & \textbf{Microsoft AutoGen} & \textbf{LangChain LangGraph} & \textbf{CrewAI} \\
\midrule
\texttt{Agent} dataclass & \texttt{ConversableAgent} & State Node definition & \texttt{Agent(role, goal, backstory)} \\
\texttt{DialogueEngine.run()} & \texttt{GroupChatManager} & StateGraph transitions & \texttt{Crew.kickoff()} loop \\
\texttt{view\_for()} mirror logic & Message router / sender roles & Conditional state router & Task message routing \\
\texttt{Budget} guardrail & \texttt{max\_consecutive\_auto\_reply} & \texttt{recursion\_limit} & \texttt{max\_iter} / rate limits \\
\texttt{manage\_context()} & \texttt{transform\_messages} filter & State reducer / window trimmer & Memory summary storage \\
\texttt{judge.py} & Benchmark evaluation agent & Evaluator node pass & Task callback validator \\
\bottomrule
\end{tabular}%
}
\end{table}

\subsection{What We Would Do Next}
If we were to continue developing this project, our next two priorities would be:
\begin{enumerate}
    \item \textbf{Structured Fact Memory:} Replace natural language summaries with a shared key-value store of confirmed alibi statements and times to avoid contradictory statements.
    \item \textbf{Improved Judge Prompts:} Refine the judge prompt with clearer rules to prevent sycophancy bias and score the inquiry based on factual contradictions rather than polite tone.
\end{enumerate}

\end{document}
